\begin{textsample}{POS Dim 1 – vem\_ed\_24 – Score 4.00 – vem\_ed\_24/t127.txt}  \label{ex:f1_pos_058}
IVES Cesu 2024: Oportunidades e desafios A \textbf{questão} “¿ Y ahora qué con COIL?
Oportunidades y desafíos” \textbf{foi} abordada por Andrea Alsina e Sofia Pekarek, da Universidad Católica de Salta (Argentina).
Com base em um relatório sobre riscos globais do Fórum Econômico Mundial, as \textbf{questões} econômicas e \textbf{ambientais} na América Latina vieram à tona, assim como suas influências nos projetos COIL, como no caso da tragédia climática recente no Rio Grande do Sul, em \textbf{que} estudantes não puderam prosseguir no Intercâmbio Virtual por falta de energia ou mesmo por estarem desabrigados.
Ambas lembraram \textbf{que} \textbf{é} preciso considerar esses aspectos externos e aqueles internos à organização, como os perfis docentes, \textbf{que} devem \textbf{ser} \textbf{identificados} e desenvolvidos em um plano de internacionalização, além de outros aspectos, como as barreiras idiomáticas, as plataformas \textbf{utilizadas}, a gestão de dados pessoais e a inclusão dos projetos COIL nos \textbf{históricos} escolares dos participantes. “COIL pressupõe criatividade e inovação para a internacionalização. \textbf{É} um \textbf{processo} do qual deriva o enriquecimento de \textbf{todas} as \textbf{partes} envolvidas”, comentou Andrea Alsina.
Sobre os perfis de professores, às vezes, um docente com dedicação exclusiva não tem disponibilidade para desenvolver projetos COIL e, por outras, um professor com “perfil táxi” (atuante em várias instituições) tem \textbf{habilidades} tecnológicas, interculturais, \textbf{linguísticas} e interesse em realizar Intercâmbios Virtuais.
Para isso \textbf{é} preciso dar incentivo para \textbf{que} os professores se envolvam em projetos COIL e incluir essa metodologia no currículo dos cursos superiores. “Os medos \textbf{são} barreiras, mas temos \textbf{que} nos despojar do erro e apostar na singeleza do projeto, entendendo \textbf{que} \textbf{todos} vamos \textbf{aprender}”, recomendou Andrea Alsina. “\textbf{É} importante \textbf{identificar} o perfil do plantel docente – cosmopolita, digital, táxi... e compreender suas forças e fraquezas”, completou Sofia Pekarek.
AGENDA – Eventos no segundo semestre de 2024 RedAES – Rede de Apoio ao Ensino Superior – 29 de agosto, on-line e gratuito.
Informações em: https://www.even3.com.br/3-jornada-redaes/ Red LatAM COIL – 4º Congresso – 4 a 6 de setembro, on-line.
Informações em: https://encr.pw/GKFNW Jornada de PCIs – 03 de outubro, on-line e gratuito (informações pelo cesu.pci@fatec.sp.gov.br) IVEC – International Virtual Exchange Conference – 21 a 24 de outubro, on-line.
Informações em: https://iveconference.org/

% matched lemmas: ambiental, aprender, habilidade, histórico, identificar, linguístico, parte, processo, que, questão, ser, todo, utilizar
\end{textsample}
